\documentclass[10pt, aspectratio=169]{beamer}

% handout
% \documentclass[handout, 9pt]{beamer}

\usepackage{clases}

\title{Investigación Operativa}
\subtitle{Método de diccionarios - $\infty$ soluciones y soluciones degeneradas}
\date{}
\author{Nazareno Faillace Mullen}
\institute{Departamento de Matemática - FCEN - UBA}
% \titlegraphic{\hfill\includegraphics[height=1.5cm]{logo.pdf}}

\begin{document}

\maketitle


\begin{frame}{Infinitas soluciones óptimas}
    \small
    Por ejemplo, considerar el siguiente diccionario para un problema de maximización:
    \[
        \begin{array}{ccccccccc}
            w_1 & = & 3 & + & x_2  & - & 2w_2 & + & 7x_3 \\
            x_1 & = & 1 & - & 5x_2 & + & 6w_2 & - & 8x_3 \\
            w_3 & = & 4 & + & 9x_2 & + & 2w_2 & - & x_3  \\ \hline
            z   & = & 8 &   &      & - & w_2  & - & x_3  \\
        \end{array}
        \Rightarrow
        (\underbrace{1}_{x_1},\underbrace{0}_{x_2},\underbrace{0}_{x_3},\underbrace{3}_{w_1},\underbrace{0}_{w_2},
        \underbrace{4}_{w_3}) \text{ es una solución óptima.}
    \]
    Si aumentamos el valor de $x_2$, la función objetivo no disminuye. Toda solución óptima debe cumplir que
    $x_3=w_2=0$ (pero no necesariamente
    $x_2=0$). Para tales soluciones, el diccionario implica que:

    \[
        \begin{array}{ccccc}
            w_1 & = & 3 & + & x_2  \\
            x_1 & = & 1 & - & 5x_2 \\
            w_3 & = & 4 & + & 9x_2 \\
        \end{array}
    \]
    Luego, cada solución óptima se obtiene asignando valores a $x_2$ tales que:
    \[  \begin{cases}
            \begin{array}{cccc}
                - & x_2  & \leq & 3 \\
                & 5x_2 & \leq & 1 \\
                - & 9x_2 & \leq & 4 \\
                \multicolumn{2}{r}{x_2} & \geq & 0
            \end{array}
            \iff x_2\in [0,\frac{1}{5}]
        \end{cases}
    \]
%    \footnotesize

\end{frame}

\begin{frame}{Infinitas soluciones óptimas}
   Es decir, podemos asignar cualquier valor en $[0,\frac{1}{5}]$ a $x_2$
    y vamos a obtener otra solución factible óptima. Por ejemplo, si tomamos $x_2=\frac{1}{6}$
    obtenemos otra solución óptima:

   $(\frac{1}{6}, \frac{1}{6}, 0, \frac{19}{6}, 0, \frac{33}{6})$ [que NO es solución básica, pero si factible]

    \vspace{1cm}
    En general, \textbf{si llegamos a una solución óptima y hay alguna variable no básica con costo reducido igual a
    $0$, entonces el problema tiene infinitas soluciones óptimas}.
\end{frame}

\begin{frame}{Problemas no acotados}

%    Recordar que la variable que sale de la base es la variable básica cuya no-negatividad impone la cota superior más
%    restrictiva al incremento de la variable que va a entrar en la base.\\
    Ejemplo (objetivo de maximizar):
    \[
        \begin{array}{ccccccccc}
            x_2 & = & 5 & + & 2x_3 & - & x_4  & - & 3x_1 \\
            x_5 & = & 7 &   &      & - & 3x_4 & - & 4x_1 \\ \hline
            z   & = & 5 & + & x_3  & - & x_4  & - & x_1  \\
        \end{array}
    \]

    La variable que va a entrar a la base es $x_3$, pero ninguna de las variables básicas $x_2$ ni $x_5$
    imponen alguna restricción sobre cuánto puede aumentar $x_3$ $\Rightarrow$
    el problema no está acotado: puedo aumentar $x_3$ tanto como yo quiera (por lo tanto, $z$
    puede aumentar tanto como yo quiera).\\
    \vskip 0.3cm
    En general, \textbf{si no existen candidatos para dejar la base, el problema es no acotado.}
\end{frame}

\begin{frame}{Soluciones degeneradas}
    Se dice que \textbf{una solución básica es degenerada si una o más de sus variables básicas valen cero}.\\
    Ejemplo:
    \[
        \begin{array}{ccccccccc}
            x_3 & = & 0.5 &   &      & - & 0.5x_4 &   &      \\
            x_5 & = &     & - & 2x_1 & + & 4x_2   & + & 3x_4 \\
            x_6 & = &     &   & x_1  & - & 3x_2   & + & 2x_4 \\ \hline
            z   & = & 4   & + & 2x_1 & - & x_2    & - & 4x_4 \\
            \multicolumn{9}{c}
            {\text{\textcolor{red}{Solución actual: $(0,0,\mathbf{\frac{1}{2}},0,\mathbf{0},\mathbf{0})$}}}
        \end{array}
    \]
    $x_1$ entrará a la base, pero observemos qué ocurre en la segunda ecuación:
    \[ 0 \leq x_5 = -2x_1 \quad \Leftrightarrow \quad x_1 \leq 0\]
    $\Rightarrow$ el crecimiento de $x_1$ está limitado a 0
    $\Rightarrow$ ni el valor de $x_1$ ni el de las otras variables va a cambiar, y el valor de $z$
    sigue siendo el mismo.\\
\end{frame}

\begin{frame}{Soluciones degeneradas}
    Aún así, debemos \textit{pivotear} para permitir que $x_1$ entre a la base y que $x_5$ salga.
    \[
        \begin{array}{ccccccccc}
            x_1 & = &     &   & 2x_2 & + & 1.5x_4 & - & 0.5x_5 \\
            x_3 & = & 0.5 &   &      & - & 0.5x_4 &   &        \\
            x_6 & = &     & - & x_2  & + & 3.5x_4 & - & 0.5x_5 \\ \hline
            z   & = & 4   & + & 3x_2 & - & x_4    & - & x_5    \\
            \multicolumn{9}{c}
            {\text{\textcolor{red}{Solución actual: $(\mathbf{0},0,\mathbf{\frac{1}{2}},0,0,\mathbf{0})$}}}
        \end{array}
    \]
    Puesto que esta iteración de SIMPLEX no produjo un cambio en el valor $z$, es una \textbf{iteración degenerada}.\\
    \vskip 0.3cm
    En el ejemplo, la próxima iteración también será degenerada, pero la posterior no. Puede ocurrir una racha de
    iteraciones degeneradas, pero en la gran mayoría de los casos, la racha es interrumpida por una iteración no
    degenerada.
\end{frame}

\begin{frame}{Soluciones degeneradas - Ciclos}
    Podría ocurrir que las soluciones degeneradas provoquen ciclos infinitos en el SIMPLEX. Ejemplo (Chvátal, 1983):
    \[
        \begin{array}{rrrrrrrrrrr}
            x_5 & = &   & - & 0.5x_1 & + & 5.5x_2 & + & 2.5x_3 & - & 9x_4  \\
            x_6 & = &   & - & 0.5x_1 & + & 1.5x_2 & + & 0.5x_3 & - & x_4   \\
            x_7 & = & 1 & - & x_1    & + &        &   &        &   &       \\ \hline
            z   & = &   &   & 10x_1  & - & 57x_2  & - & 9x_3   & - & 24x_4 \\
        \end{array}
    \]
    Luego de la primera iteración ($x_1$ entra, $x_5$ sale):
    \[
        \begin{array}{rrrrrrrrrrr}
            x_1 & = &   &   & 11x_2 & + & 5x_3  & - & 18x_4  & - & 2x_5  \\
            x_6 & = &   & - & 4x_2  & - & 2x_3  & + & 8x_4   & + & x_5   \\
            x_7 & = & 1 & - & 11x_2 & - & 5x_3  & + & 18x_4  & + & 2x_5  \\ \hline
            z   & = &   &   & 53x_2 & + & 41x_3 & - & 204x_4 & - & 20x_5 \\
        \end{array}
    \]
    \large
    \[ \mathbf{\vdots} \]
\end{frame}

\begin{frame}{Soluciones degeneradas}
    \[\mathbf{\vdots}\]
    Luego de la quinta iteración:
    \[
        \begin{array}{rrrrrrrrrrr}
            x_5 & = &   &   & 9x_6  & + & 4x_1   & - & 8x_2   & - & 2x_3   \\
            x_4 & = &   & - & x_6   & - & 0.5x_1 & + & 1.5x_2 & + & 0.5x_3 \\
            x_7 & = & 1 &   &       & - & x_1    &   &        &   &        \\ \hline
            z   & = &   &   & 24x_6 & + & 22x_1  & - & 93x_2  & - & 21x_3  \\
        \end{array}
    \]

    Luego de la sexta iteración:
    \[
        \begin{array}{rrrrrrrrrrr}
            x_5 & = &   & - & 0.5x_1 & + & 5.5x_2 & + & 2.5x_3 & - & 9x_4  \\
            x_6 & = &   & - & 0.5x_1 & + & 1.5x_2 & + & 0.5x_3 & - & x_4   \\
            x_7 & = & 1 & - & x_1    & + &        &   &        &   &       \\ \hline
            z   & = &   &   & 10x_1  & - & 57x_2  & - & 9x_3   & - & 24x_4 \\
        \end{array}
    \]
    ¡Volvimos al mismo diccionario del principio!
\end{frame}

\begin{frame}{Soluciones degeneradas - Regla de Bland}
    \begin{tcolorbox}[colback=blue!25!white, colframe=blue!45!white]
        \textbf{Regla de Bland}:
        \begin{enumerate}
            \item La variable entrante será la variable no básica \textbf{de menor índice}
            con coeficiente positivo en la f.o.
            \item
            La variable que sale se elige como siempre. Si hay empates entre variables, sale la que tenga menor índice.
        \end{enumerate}
    \end{tcolorbox}
    \pause
    \vskip 0.5cm
    Se puede demostrar que, utilizando la regla de Bland, SIMPLEX no cicla. [Demo: \textit{Linear Programming}
    , Chvátal, 1983, p. 37]\\
    \vskip 0.2cm
    \textbf{Obs:}
    no es necesario utilizar la regla de Bland en cada iteración. Podríamos comenzar a emplearla, por ejemplo, luego de
    algunas iteraciones degeneradas y descartarla cuando se produzca una iteración no degenerada.
\end{frame}

\begin{frame}{Soluciones degeneradas - Regla de Bland}
    En el ejemplo donde SIMPLEX ciclaba, veamos cómo la regla de Bland nos ayuda a salir del ciclo. En la quinta
    iteración teníamos el siguiente diccionario:
    \[
        \begin{array}{rrrrrrrrrrr}
            x_5 & = &   &   & 9x_6  & + & 4x_1   & - & 8x_2   & - & 2x_3   \\
            x_4 & = &   & - & x_6   & - & 0.5x_1 & + & 1.5x_2 & + & 0.5x_3 \\
            x_7 & = & 1 &   &       & - & x_1    &   &        &   &        \\ \hline
            z   & = &   &   & 24x_6 & + & 22x_1  & - & 93x_2  & - & 21x_3  \\
        \end{array}
    \]
    Vimos que si usamos el criterio usual, $x_6$ entra a la base y así volvemos al diccionario inicial. \\
    En cambio, si aplicamos la Regla de Bland, $x_6$ no entrará a la base, sino que $x_1$
    entrará. Cada ecuación nos impone las siguientes restricciones, respectivamente:
    \[ x_1<\infty \qquad x_1\leq 0 \qquad x_1\leq 1 \]
    Luego, $x_4$ sale de la base. Si bien esta es una iteración degenerada, veamos cómo queda el nuevo diccionario.
\end{frame}

\begin{frame}{Soluciones degeneradas - Regla de Bland}
    \small
    Una vez que pivoteamos, el nuevo diccionario queda así:
    \[
        \begin{array}{rrrrrrrrrrr}
            x_1 & = &   & - & 2x_6  & + & 3x_2  & + & x_3  & - & 2x_4  \\
            x_5 & = &   &   & x_6   & + & 4x_2  & + & 2x_3 & - & 8x_4  \\
            x_7 & = & 1 & + & 2x_6  & - & 3x_2  & - & x_3  & + & 2x_4  \\ \hline
            z   & = &   & - & 20x_6 & - & 27x_2 & + & x_3  & - & 44x_4 \\
        \end{array}
    \]
    \pause
    Primero, no volvimos al diccionario inicial. Por otro lado, veamos qué sucede con la próxima iteración: $x_3$
    entrará a la base y $x_7$ saldrá de la base. Luego de efectuar el pivote, el nuevo diccionario queda:
    \[
        \begin{array}{rrrrrrrrrrr}
            x_3 & = & 1 & + & 2x_6  & - & 3x_2  & + & 2x_4  & - & x_7  \\
            x_1 & = & 1 &   &       &   &       &   &       & - & x_7  \\
            x_5 & = & 2 & + & 5x_6  & - & 2x_2  & - & 4x_4  & - & 2x_7 \\ \hline
            z   & = & 1 & - & 18x_6 & - & 30x_2 & - & 42x_4 & - & x_7  \\
        \end{array}
    \]
    Como $z$
    cambió su valor (pasó de valer 0 a valer 1), la última iteración no es degenerada. Más aún llegamos a una solución
    óptima: (1,0,1,0,1,0,0)
    \vskip 0.3cm
\end{frame}

\begin{frame}{Soluciones degeneradas - Regla de Bland}
    \textbf{En síntesis:}
    \begin{itemize}
        \item en general es útil \textbf{el criterio usual}
        para decidir qué variable entra a la base. En el caso de maximizar \textquotedblleft
        Entra a la base la variable con mayor coeficiente positivo en la f.o." (análogo para minimizar)
        \item si se viene dando una \textbf{seguidilla de iteraciones degeneradas, usar la Regla de Bland}.
    \end{itemize}
\end{frame}

\end{document}


